\documentclass[11pt,a4paper]{article}

\usepackage[T1]{fontenc}
\usepackage[utf8]{inputenc}
\usepackage{textalpha}
\usepackage[spanish]{babel}
\usepackage[a4paper,margin=2.8cm]{geometry}
\usepackage{microtype}
\usepackage{xurl}
\usepackage[hidelinks]{hyperref}

\hypersetup{
  pdftitle={Todo lo que merece lágrimas. De concursos literarios, detectores de IA y caza de brujas},
  pdfauthor={Juan José Gómez Cadenas}
}

\setlength{\parindent}{1.25em}
\setlength{\parskip}{0.35em}
\raggedbottom

\title{Todo lo que merece lágrimas\\[0.4em]\large De concursos literarios, detectores de IA y caza de brujas}
\author{Juan José Gómez Cadenas}
\date{}

\begin{document}

\maketitle

Estamos en el año 20~a.~C., en la reunión final del \emph{Praemium Stellae Errantis} ---conocido por el vulgo como \emph{Praemium Planetarium}---. Se han recibido numerosos textos ---el premio es muy codiciado, ya que paga nada menos que cien mil sestercios---. El trabajo favorito ---cuyo título final y autor se mantienen en el anonimato y solo se conocen como E. y V.--- está organizado en 12 cantos y escrito en pulcros hexámetros latinos. La historia narra las peripecias de un héroe que busca su hogar, tras combatir en una terrible guerra. Nadie duda de que V. se va a llevar el galardón ---y la pila de sestercios---, hasta que, después de la comilona en casa del presidente del jurado, Quinto Publio ---uno de los miembros externos de la comisión--- sale por peteneras.

Antes de adentrarnos en su discurso, conviene aclarar un par de datos sobre los concursos literarios en Roma. Hubo un tiempo en el que se hicieron muy célebres, en particular el \emph{Praemium Planetarium} que nos ocupa, seguido de otros casi igual de notables como el \emph{Praemium Natalis} y el \emph{Praemium Alfaguarianum}. Se hicieron, de hecho, tan populares que los patricios literarios de la época se dieron cuenta de que se les podía sacar tajada comercial con una pizca de imaginación. No es que les diera por hacer trampas, faltaría más. Si Roma no paga traidores, mucho menos va a pagar plumíferos. Claro que si en el jurado se designan unos pocos \emph{gazetarii} afines a la casa que reúnen la mayoría de los votos y unos pocos miembros externos que dan prestigio y generalmente no se enteran de nada ---mucho menos después de agasajarlos con unos cuantos \emph{convivia}, bien regados por vino de Falerno---, ¿se incumple alguna regla? Ninguna, y \emph{negotium est negotium}.

Lo malo es que el bueno de Quinto ha empinado el codo un poco de más y un poco de menos en el \emph{convivium}.

Un poco de más para recordar que está en el jurado por ser sobrino del editor y no por sus discutibles méritos literarios.
Quinto solo ha publicado un libro en su vida, \emph{De apibus quae mel non faciunt}, dedicado al estudio de las abejas que no hacen miel, y podría decirse que su éxito ha sido muy modesto.

Un poco de menos, porque, generalmente, nuestro literato se bebe hasta el vinagre de lavar las copas en los banquetes y cuando llega la reunión del jurado se limita a dormitar, soñando con sus himenópteros. Pero una molesta acidez estomacal ha contenido algo sus ansias y el resultado es un súbito ataque de imprudencia complicado con un subidón de verborrea. Así que al llegar su turno, Quinto toma la palabra (más bien la coge por el cuello) y pronuncia este discurso:

\begin{quotation}
¿Hasta cuándo, V., abusarás de la paciencia del jurado? ¿Cuánto tiempo más vamos a sufrir doce libros de una obra que ya conocíamos en veinticuatro? ¡Oh tiempos, oh certámenes!

Nobles patricios, ¡basta con leer los primeros versos! \emph{Arma virumque cano}.

¿No les suena a nada? ¿No les recuerdan al divinal Homero?: \emph{Ἄνδρα μοι ἔννεπε}.

\emph{Ἄνδρα μοι ἔννεπε}. El Poeta canta al varón, a \emph{Ἄνδρα}, en el primer verso. Y el desvergonzado V. hace lo mismo, cantando al varón y a las armas.

Estamos ante un caso de horrible plagio. Lean con atención. Su obra no es otra cosa que la \emph{Odisea}. Ulises parte de Troya y surca con gran peligro los mares. Encuentra cíclopes y hechiceras, pasa entre Escila y Caribdis, escapa de lotófagos y lestrigones, para llegar a la patria, Ítaca, y reunirse con su fiel esposa, Penélope.

¿Y qué hace E., el protagonista de este trabajo? Escapa de Troya, su hogar, recorre exactamente el mismo camino del héroe, sin escatimar cíclopes, ni Escila, ni la isla de la peligrosa Circe. Abandona a Dido, que no es sino un trasunto de Calipso, y llega por fin a Italia ---obviamente, a Ítaca---, donde le espera la bella Lavinia, similar a Penélope, con quien se unirá felizmente, no sin antes aniquilar a Turno y sus rútulos, es decir, a los pretendientes.

¡Plagio vulgar! Y lo más vulgar de todo, lo más intolerable, es el \emph{descensus Averno} del protagonista, el viaje al Hades que no es sino una descarada copia de la \emph{Νέκυια} homérica.

Patricios, miembros del jurado. Sostengo que V. ha generado su obra, E., usando IA, un \emph{Instrumentum Amanuensium}, o si lo prefieren un \emph{Ingenium Alienum}. Sé de buena tinta que el ciudadano Ático emplea a un ejército de \emph{librarii} y con un ejército de esclavos copistas y otro de \emph{anagnostae} para leer y corregir esas copias, se puede producir de manera totalmente artificial el texto que nos ocupa.
\end{quotation}

El jurado se queda estupefacto. Todos, menos el poeta Horacio, que ni se inmuta, dándole sorbitos a su Falerno, con una sonrisa irónica en los labios. Quinto Publio se viene arriba y se salta las reglas de mantener en el anonimato el nombre del escritor y de su obra:

---¡Exijo que la \emph{Eneida} de Virgilio sea descalificada como un plagio IA de la \emph{Odisea}!

Hace un par de días, mi admirado Kiko Llaneras levantaba una considerable polvareda en el erial patrio, tras destapar que la novela \emph{Donde el agua trabaja}, del autor David Silva Rojas, tiene un 99~\% de texto escrito con IA, según el detector Pangram. La obra en cuestión había ganado los nada desdeñables 20.000 euros del Premio de Novela Breve Almadía-Ventosa Arrufat, un premio cuyas bases, por cierto, prohibían explícitamente usar IA para escribir los textos.

En su newsletter, Kiko se muestra prudente. Empieza tirando balones fuera en lo que se refiere a la fiabilidad de Pangram, asegurando que en fragmentos cortos se le pueden llevar al huerto, cosa que ya ha demostrado en esta revista el astuto Ángel Fernández. Pero, asegura Kiko, en textos largos es «mucho más estable». Se queda corto. Pangram puede chirriar en un texto corto (muy corto) y está diseñado para minimizar los falsos positivos. Esto quiere decir que si un autor tiene la paciencia y la maña de reescribir un borrador generado con IA introduciendo bastante de su cosecha como para que el detector se salga de «la zona IA» de su mapa de parámetros, el bendito de Pangram lo va a dar por bueno. Cuando Kiko dice que Pangram no es infalible, se salta una precisión muy importante: \emph{No es infalible porque les concede a los humanos el beneficio de la duda}. Así que es posible generar textos con IA, maquillarlos lo suficiente y darle a Pangram humano por máquina. \emph{Pero en un texto largo las posibilidades de que Pangram dé, falsamente, 99~\% positivo son astronómicamente bajas, cero a todos los efectos}. De lo que se pueden sacar varias conclusiones.

\begin{enumerate}
\item El autor escribió la obra usando IA.
\item El autor lo niega (según entiendo). Si es el caso, miente.
\item Un concurso literario que prohíbe el uso de la IA debería molestarse en pasar uno o varios detectores de texto fiables para detectar a los tramposos. No hacerlo es falta de seriedad, como mínimo.
\item El jurado del premio estaba en las berzas.
\end{enumerate}

Todo muy edificante. No es la primera vez que pasa, además. El escándalo del Goncourt precedía nuestra pifia patria\footnote{\url{https://www.lemonde.fr/en/pixels/article/2026/09/25/french-literary-prize-removes-novel-from-longlist-after-claims-of-using-ai_6757946_13.html}}, aunque, eso sí, en el certamen francés pillaron al tramposo antes de dar el premio y lo sacaron de la lista.

Dice Kiko Llaneras que no quiere una caza de brujas. Alega muy píamente que él usa la IA para todo ---pero servidor, que es un malpensado, ha pasado su texto por Pangram y sale, como era de esperar, 100~\% humano---. Luego ofrece no menos píos escenarios de colaboración hombre-máquina.

Ya he dicho que admiro a Llaneras. Es un tipo que sabe de lo que habla y habla de lo que sabe. Sus newsletters son de lo mejor que puede encontrarse en \emph{El País}. Es ameno, amable, escribe bien y sabe hacer sus cuentas. Pero lo cierto es que discrepo de él ---y de mi no menos admirado amigo Ángel--- en eso de la caza de brujas.

Yo estoy, que dirían los de la Mandrágora, por la hoguera.

En mis últimos artículos he venido desarrollando un argumento que aquí vuelve a repetirse: \emph{el problema no es la IA, somos nosotros}. Los ataques informáticos, la pornografía, las falsas cuentas en redes sociales o los timos online son pecados humanos. La IA no hace sino amplificarlos.

El plagio es uno de los pecados literarios más mortales, posiblemente el más grave de todos, peor aún, y ya es decir, que la prosa pedante. Y decirle a la IA que te escriba un texto es una forma de plagio a lo grande, porque la IA se va a limitar a sintetizar, siguiendo sus prompts y sus pesos, un derivado de las obras de las que se ha alimentado.

¿Y no es eso lo que hacemos siempre? ¿No es eso el proceso de escritura? Después de todo, ¿no es verdad que Virgilio plagió a Homero?

Volvamos a la tarde de hace dos milenios y pico. Tras la acusación de Quinto Publio, los miembros del jurado se miran consternados, sin saber por dónde salir. Horacio se levanta entonces de su \emph{triclinium} y lee los siguientes versos:

\begin{verse}
\emph{Ibant obscuri sola sub nocte per umbram}
\end{verse}

Se hace un silencio. Quinto, que es un poco espeso, no entiende qué tiene que ver la soledad de la noche y la oscuridad con la que caminan Eneas y la Sibila por las vacías moradas de Dite y sus reinos sin sustancia. Horacio se lo tiene que explicar, como quien le repasa la lección a un escolar algo lerdo.

---Dime qué nos cuenta el poeta, Quinto.

---Está muy claro ---dice el acusador---. El héroe y su compañera caminan solitarios por la oscuridad de la noche, entre las sombras.

---No, no dice eso ---corrige Horacio---. Eso es lo que habría escrito uno de los \emph{Instrumenta Amanuensium} ---ya sabes, los IA--- de la casa de Ático. Virgilio escribe algo muy diferente. Dice: \emph{Iban oscuros por la soledad de la noche}.

¿Por qué es tan diferente? ¿Qué ha entendido Horacio que se le escapa al buen Quinto? El grandioso Jorge Luis Borges lo explicaría milenios más tarde\footnote{\url{https://lectoresaudaces.blogspot.com/2011/09/mas-cosas-sobre-la-eneida-de-borges-el.html}}:

\begin{quotation}
Virgilio no nos dice que los aqueos aprovecharon los intervalos de la oscuridad para entrar en Troya, habla de los amistosos silencios de la luna. No escribe que Troya fue destruida, escribe “Troya fue”. No escribe que un destino fue desdichado, escribe “De otra manera lo entendieron los dioses”. Para expresar lo que ahora se llama panteísmo nos deja estas palabras: “Todas las cosas están llenas de Júpiter”. Virgilio no condena la locura bélica de los hombres, dice “El amor del hierro”. No nos cuenta que Eneas y la Sibila erraban solitarios bajo la oscura noche entre sombras, escribe:

\begin{verse}
\emph{Ibant obscuri sola sub nocte per umbram}
\end{verse}
\end{quotation}

Virgilio introduce la hipálage más célebre de la historia de la literatura. Andar solitarios en la noche oscura es una pura descripción de la realidad. Andar oscuros en la noche solitaria es poesía, como es poesía la elipsis que nombra los amistosos silencios de la luna, como es poesía esa otra maravillosa línea:

\begin{verse}
\emph{Sunt lacrimae rerum}.
\end{verse}

Tras la caída de Troya, Eneas y su gente llegan a Cartago, donde les recibe amistosamente la malograda reina Dido. Extrañamente, en las paredes de un templo, encuentran pinturas que representan sus desgracias. Allí están las imágenes de Príamo, de Aquiles, de Héctor, del propio Eneas. Entiende el héroe, en ese momento, que aquella tierra extraña que les recibe no es, sin embargo, ajena a sus pesares. Y es en este contexto donde Virgilio nos ofrece el maravilloso verso:

\begin{verse}
\emph{Sunt lacrimae rerum, et mentem mortalia tangunt}
\end{verse}

La traducción literal es muy extraña. Sería algo así como «Hay lágrimas de las cosas y las cuitas mortales tocan a la mente». Y ahí está la magia, escondida a plena luz. Estas «lágrimas de las cosas» podrían referirse a que las cosas lloran. El poeta estaría humanizando los objetos del mundo, haciendo que el orbe llorara por la desgracia de los troyanos. Pero otra posible traducción ---la que yo prefiero--- sería «todo lo que merece lágrimas» (ver el excelente texto\footnote{\url{https://clasicos.hypotheses.org/2099}}). Todo lo que merece lágrimas universaliza nuestra compasión, nos hace llorar con Dido por las penurias de Eneas y con Virgilio por el destino de Dido, nos hace lamentarnos ante todas y cada una de las desgracias humanas, las que conocemos y las que ignoramos, las que nos pillan de lejos y de cerca, las que nos laceran y las que pasan de largo. Llorar por todo lo que merece lágrimas es, quizás, la definición más profunda y más redentora de un ser humano.

Cuando Horacio acaba su exposición, esas lágrimas se asoman a los ojos de los patricios que le rodean. Porque los versos les han traído a la mente la destrucción de Troya y la de Cartago, la muerte de Dido y la de Julieta ---sin saberlo, lloran por la adolescente enamorada que imaginaría Shakespeare siglos más tarde---. Lloran por la muerte de Anquises, que muere otra vez cuando perdemos a nuestro propio padre. Lloran, con el corazón atravesado de remordimiento por el asesinato de Astianacte, arrojado desde la alta muralla, y lloran las mismas lágrimas de alivio porque Ascanio vive y con él la esperanza de que cada uno de sus hijos ---cada uno de nuestros hijos--- nos supere en longevidad y fortuna, en felicidad y suerte. Lloran por todo lo que merece lágrimas, la soledad de la noche y la oscuridad de los hombres y los vacíos reinos del mundo de las sombras que aguardan a todos los humanos, aunque no falta mucho para que cierto nazareno asegure que el mundo de las sombras no existe y el destino de las almas es el Paraíso. Lloran como llora Marty, cuando Rust le dice:

\begin{quotation}
«Antes solo había oscuridad. Si me preguntas, creo que la luz va ganando.»
\end{quotation}

Todo lo que merece lágrimas, todo lo humano, merece ser escrito por humanos. La razón por la que Virgilio no plagia la \emph{Odisea} es que su obra, que, a propósito, sigue el patrón homérico, añade a la literatura universal un héroe consumido por el deber en lugar del astuto y desalmado Ulises, una reina trágica cuyo aciago destino se va a repetir una y otra vez en los siguientes milenios, una muerte airada que degrada al protagonista en la última línea, un paseo oscuro en la soledad del Hades y todas las cosas que merecen lágrimas. Virgilio no plagia porque crea donde no había, saca, literalmente, luz de la oscuridad. Como tampoco plagia Dante a Virgilio ---que le lleva de la mano a los infiernos--- ni Joyce a ninguno de ellos. La literatura es un río del que beben todas las almas, las de los vivos y las de los muertos.

Creo que cuando un concurso literario pide que no se use IA, lo mínimo que puede hacer es pasar un detector fiable. Pretender que Pangram se equivoca cuando da 99~\% rojo (en un texto de más de unos cientos de palabras, no digamos ya en 300 páginas de novela) es mentir sin escrúpulos. Lamentarse de que los policías de la IA queremos cancelar a los creadores culturales es hipócrita. No seré yo quien tire la primera piedra si alguien firma sus artículos como Claude y Pepe Luis, si añade en un aviso que el suyo es un texto cogenerado con IA. Otra cosa es ocultar (la coautoría de la IA) o mentir (negándola). Seamos serios. Nadie quiere cancelar a nadie. Solo se pide un poco de honestidad.

No afirmo que un texto generado o cogenerado por IA carezca de interés. Tengo muchos ejemplos de lo contrario. Sostengo, sin embargo, firme, apasionadamente, que las lágrimas de las cosas no están en ellos.




\end{document}
